\chapter{Conclusão}
\label{cap:conclusao}

%% LC (resolvido 12/11/2026): resposta explícita à questão de pesquisa e ligação
%% de cada objetivo à sua evidência.
A questão que orientou esta pesquisa foi se é viável aplicar \textit{Bounded
Model Checking} a um \textit{firmware} de controle de voo de porte industrial,
com tempo de verificação por módulo compatível com o uso prático e produzindo
evidências auditáveis sobre propriedades de memória e de aritmética. A resposta
é afirmativa, com um alcance que precisa ser dito junto: a viabilidade foi
demonstrada por módulo, sobre \textit{harnesses} que isolam cada unidade, e não
para o sistema completo. Vinte e dois módulos do PX4 Autopilot foram
verificados em duas campanhas; cada verificação da primeira concluiu em menos
de noventa segundos, e todo veredito pode ser refeito a partir dos artefatos
publicados.

Os três objetivos específicos foram atendidos da seguinte forma. O objetivo~O1,
propor a metodologia, corresponde às oito etapas do
Capítulo~\ref{cap:metodologia}. O objetivo~O2, aplicá-la e classificar os
candidatos, corresponde às duas campanhas do Capítulo~\ref{cap:experimentos}:
três contraexemplos triados na primeira e dez módulos com desfecho individual
na segunda (Tabelas~\ref{tab:resultados} e~\ref{tab:desfechos-segunda}). O
objetivo~O3, avaliar a viabilidade, corresponde aos tempos e às configurações
por módulo da Tabela~\ref{tab:resultados} e às respostas às questões
experimentais.

\section{Síntese dos Resultados}
\label{sec:conclusao-sintese}

%% LC (resolvido 12/11/2026): síntese condensada em um parágrafo por achado.
\textbf{Dez módulos sem violação.} Na primeira campanha, os
\textit{harnesses} de dez dos doze módulos foram verificados sem violação das
propriedades de segurança de memória e de \textit{overflow} habilitadas. A
reexecução com as asserções de desenrolamento ligadas mostrou que nenhum laço
desses \textit{harnesses} excede o limite usado, de modo que o resultado não
depende de truncamento silencioso.

\textbf{Temperatura do BMI088.} O verificador produziu uma leitura de
$86{,}25\,^{\circ}$C, acima do limite operacional de $85\,^{\circ}$C do
\textit{datasheet}. Os mantenedores recusaram a observação como defeito, e ela
permanece como especificação disputada.

\textbf{Contador do FIFO do BMI088.} O contraexemplo, com
\texttt{count}~$= 9088$ para um FIFO de $1.024$~bytes, era um falso positivo: o
driver de produção contém três guardas que o \textit{harness} omitia.

\textbf{GPS uBlox.} O verificador encontrou, no \textit{harness}, uma escrita
fora dos limites em \texttt{dumpGpsData()}. Uma prova indutiva sobre o laço real
do driver mostrou que o estado violador é inalcançável na produção, e um ensaio
em bancada foi consistente com essa conclusão. A guarda proposta foi enquadrada
pelo mantenedor como endurecimento defensivo.

\textbf{Segunda campanha.} Em sete de dez módulos complementares, o efeito
previsto foi reproduzido executando o código real do PX4 sob sanitizadores.
Três desses defeitos foram corrigidos pelo próprio projeto, de forma
independente; dois foram reportados neste trabalho (RoboClaw e batmon) e
aguardavam triagem no fechamento do texto; um já havia sido relatado por
terceiro; e um não foi reportado, pelo impacto desprezível. Um falso positivo
foi barrado pela conferência antes de qualquer relato.

\textbf{Validação em bancada.} Numa Pixhawk~6X~RT, o driver do RoboClaw
transmitiu o byte de dados 130, fora da faixa de 0 a 127 do protocolo, e deixou
de fazê-lo com a correção aplicada. O controlador de motores foi emulado.

\textbf{Correção na ferramenta.} A campanha expôs um defeito no modelo de
\texttt{snprintf} do ESBMC, confirmado e corrigido pelos mantenedores.

\section{Contribuições}
\label{sec:conclusao-contribuicoes}

%% LC (resolvido 12/11/2026): lista alinhada às cinco contribuições da
%% Introdução, com a numeração (i) a (viii) das etapas.
As contribuições são as cinco enunciadas na Introdução
(Seção~\ref{sec:contribuicoes}).

\begin{enumerate}
\item \textbf{Metodologia em oito etapas.} Seleção de módulos~(i),
instrumentação com \textit{stubs}~(ii), especificação de propriedades~(iii),
execução do verificador~(iv), análise de contraexemplos com conferência contra
o código de produção~(v), reporte~(vi), validação cruzada e verificação de
\textit{patch}~(vii), e reprodução em outro autopiloto com integração
contínua~(viii).
\item \textbf{Evidência experimental delimitada pelo modelo.} Dez módulos sem
violação até o limite de desenrolamento, com tempos por módulo medidos em três
execuções.
\item \textbf{Triagem documentada.} Três contraexemplos com três desfechos
distintos, e a prova formal de que o estado violador do GPS não ocorre no
driver de produção.
\item \textbf{Segunda campanha.} Dez módulos verificados sobre o código
original, com defeitos reproduzidos por \textit{replay}, um deles observado em
bancada, dois relatos ao PX4 e uma correção aceita no ESBMC.
\item \textbf{Artefato público.} \textit{Harnesses}, registros, programas de
\textit{replay}, capturas de bancada e roteiros de reprodução.
\end{enumerate}

A ausência de teste unitário direto para \texttt{dumpGpsData()} foi conferida
por busca na árvore do PX4 (Seção~\ref{sec:caso-gps}). Ela evidencia falta de
cobertura unitária direta, e não de exercício indireto por simulação ou por
testes de integração.

\section{Limitações do Estudo}
\label{sec:conclusao-limitacoes}

\textbf{Cobertura.} Os arquivos-alvo da primeira campanha somam cerca de
$1{,}5\,\%$ das $710$~mil linhas de C/C++ de \texttt{src/}. Ficaram de fora os
mecanismos de \textit{failsafe} do Commander, a lógica de \textit{geofence}, os
barramentos CAN/UAVCAN, a ponte RTPS/uXRCE-DDS e os drivers de fluxo óptico e
telêmetro, as mesmas exclusões da Seção~\ref{sec:ameacas-validade}.

\textbf{Limite de desenrolamento.} As garantias valem até o limite $k$
adotado: 10 como padrão, 15 no Navigator/Mission e 300 no MAVLink Parser
(Seção~\ref{subsec:escolha-k}). A campanha original usou
\texttt{--no-unwinding-assertions}; a reexecução com as asserções ligadas
afastou o risco de truncamento nos dez módulos sem violação, mas, no GPS, o
laço do \texttt{memcpy} excede $k = 10$ e não é coberto por inteiro.

\textbf{Modelo.} Os \textit{stubs} não representam temporização, atrasos de
barramento nem falhas intermitentes, e a análise modular não captura defeitos
que surgem da interação entre tarefas concorrentes.

%% LC (resolvido 12/11/2026): limitação sobre defeitos confirmados.
\textbf{Eficácia em encontrar defeitos.} A primeira campanha não produziu
defeito confirmado algum: de três contraexemplos, um era falso positivo, um foi
recusado e um é inalcançável na produção. A segunda campanha reproduziu
defeitos no código real, e três deles coincidem com correções feitas pelo
próprio PX4; ainda assim, nenhum relato deste trabalho havia sido aceito pelos
mantenedores no fechamento do texto. A eficácia do método em encontrar defeitos
que o projeto reconheça como tais está, portanto, demonstrada de forma
indireta.

%% LC (resolvido 12/11/2026): limitação sobre esforço humano.
\textbf{Esforço humano.} O tempo de máquina é pequeno, mas o custo do método é
dominado pelo trabalho de escrever \textit{harnesses}, triar contraexemplos e
conferir o código de produção. Esse esforço não foi medido, e nenhuma afirmação
sobre custo total de adoção se sustenta nos dados deste trabalho.

\textbf{Validação em bancada.} Os ensaios cobrem dois achados. No do RoboClaw,
o controlador de motores foi emulado, o \textit{firmware} precisou de duas
alterações declaradas e só um dos dois defeitos foi exercitado; no do GPS, não
havia receptor conectado.

%% LC (resolvido 12/11/2026): aplicabilidade ao SAGRES.
\textbf{Relação com o projeto SAGRES.} O projeto motivou a pesquisa, mas não
foi levantada a configuração de aeronave que ele utiliza, nem quais dos módulos
verificados fazem parte dela. A aplicabilidade dos resultados às suas missões é
uma hipótese, e não um resultado.

\section{Trabalhos Futuros}
\label{sec:conclusao-futuros}

\subsection{Expansão da cobertura}

%% LC (resolvido 12/11/2026): Navigator retirado da lista de não analisados;
%% LPE definido; citações mantidas só onde sustentam a frase.
A metodologia pode ser estendida aos subsistemas não analisados: os mecanismos
de \textit{failsafe} do Commander, a lógica de \textit{geofence}, os barramentos
CAN/UAVCAN e o estimador de posição local do PX4 (\textit{Local Position
Estimator}, LPE). Nos módulos já verificados, cabe ampliar os cenários; no
MAVLink Parser, por exemplo, sequências de mensagens malformadas do tipo
descrito por Ficco et al.~\cite{Ficco2022}.

\subsection{Validação em bancada e em simulação}
\label{sec:trabalhos-futuros-hw}

%% LC (resolvido 12/11/2026): seção numerada, para que as referências apontem
%% para ela.
A validação em bancada pode ser estendida aos demais achados da segunda
campanha. O sensor de vento FT7 e o batmon são os candidatos imediatos: o
primeiro usa a mesma interface serial do ensaio do RoboClaw, e o segundo
exigiria emular uma bateria SMBus. Para o RoboClaw, falta observar a reação de
um controlador real ao byte fora da faixa, que o manual não especifica. Para o
GPS, falta o sentido do receptor para a placa.

Os ensaios feitos injetam a condição de disparo diretamente, por parâmetro e
por comando de console. Eles mostram que o defeito ocorre na placa quando a
condição se apresenta, mas não que a cadeia de módulos a montante a produz em
operação. Responder a essa pergunta exige exercitar a pilha completa, e a
simulação \textit{software-in-the-loop} é o ambiente adequado, sem risco ao
equipamento.

\subsection{Integração em \textit{pipeline} de desenvolvimento}

Um fluxo do GitHub Actions já executa os \textit{harnesses} do IMU e do GPS a
cada \textit{push} (Seção~\ref{subsec:ardupilot-ci}). Os passos seguintes são a
verificação incremental, restrita aos módulos modificados, a medição do tempo
de relógio sob configuração controlada e a proposição do fluxo à integração
contínua oficial do projeto.

\subsection{Verificação composicional e concorrente}

A verificação módulo a módulo não detecta defeitos de composição. Dois recursos
do ESBMC atacam essa limitação. Os contratos de função permitem substituir um
\textit{stub} não determinístico por uma especificação verificada do que a
função exige e garante. A verificação de programas concorrentes, não explorada
aqui, permitiria tratar condições de corrida entre as tarefas que o PX4 executa
sob o escalonador do NuttX; exigiria modelar o escalonamento e a semântica de
publicação do uORB.

\subsection{Propriedades além do software}

As propriedades verificadas referem-se a estados do programa. O que uma
violação significa para o veículo, isto é, em que condição de voo ocorreria e
com que consequência, exige acoplar à verificação um modelo abstrato da
dinâmica do sistema físico. Essa foi uma pergunta da banca de qualificação que
este trabalho não responde.

\section{Considerações Finais}
\label{sec:conclusao-finais}

As três questões experimentais tiveram as seguintes respostas. Quanto à
\textbf{QE1}, todo módulo da primeira campanha foi verificado em menos de
noventa segundos em \textit{hardware} de prateleira, o que demonstra
viabilidade por módulo. Quanto à \textbf{QE2}, a resposta é parcial: o BMC
expôs, no \textit{harness} do GPS, uma violação em função sem teste unitário
direto, mas o estado é inalcançável na produção, e nenhum defeito do driver foi
exposto na primeira campanha. Quanto à \textbf{QE3}, nove dos doze módulos
foram verificados com a invocação padrão, e três exigiram opções específicas.

%% LC (resolvido 12/11/2026): alegações de custo e de casos de canto qualificadas.
Os resultados indicam que o BMC é aplicável a módulos de \textit{firmware} de
VANTs em produção com tempo de máquina pequeno. Não permitem afirmar que o
custo total seja compatível com um ciclo de desenvolvimento, porque o esforço
humano e o tempo em integração contínua não foram medidos. Também não se
sustenta, para a primeira campanha, a afirmação de que a abordagem alcança
casos extremos que os testes não cobrem: o único caso extremo encontrado ali
não ocorre no driver real.

%% LC (resolvido 12/11/2026): argumento final apoiado na segunda campanha.
O valor da abordagem aparece com mais clareza na segunda campanha, e por uma
razão metodológica. Quando a verificação foi feita sobre o código original, e
cada candidato só foi aceito depois de reproduzido fora do verificador, sete de
dez módulos exibiram defeitos que o código real de fato tem, três deles
corrigidos pelo próprio projeto sem conhecimento deste trabalho. A primeira
campanha mostrou o que acontece sem essa disciplina: contraexemplos formalmente
válidos para defeitos que não existem.

O que este trabalho entrega é, assim, uma metodologia de verificação formal
para \textit{firmware} de sistemas críticos, aplicada ao PX4 como estudo de
caso. Os achados são produto da aplicação do método, e não a sua medida:
aplicado a outro \textit{firmware}, ele poderia não revelar violação alguma e
ainda assim produzir evidência da mesma natureza, limitada, reproduzível e
auditável. A lição que atravessa as duas campanhas é que um contraexemplo é o
início da triagem, e que a conferência contra o código de produção é a etapa
que separa um achado de um artefato do modelo.
